\section{Bernoulli の微分方程式}
特に非線形常微分方程式は多くの場合、解くことができない。

一階線形常微分方程式の左辺に対して
右辺が$Q(x) y^n \: (n\neq 0,\,1)$になっているものを Bernoulli の微分方程式
% 、$f(x)y^2 + g(x)y + h(x)$になっているものをリッカチの微分方程式
という。
\begin{gather}
  \frac{dy}{dx} + P(x)y = Q(x)y^n \label{eq:appendix-ODE-Bernoulli}
  % \frac{dy}{dx} + P(x)y = f(x)y^2 + g(x)y + h(x) \label{eq:appendix-ODE-Riccati}\\
\end{gather}

Bernoulli の微分方程式は両辺を$y^n$で割って、
\begin{gather}
  \frac{1}{y^n} \frac{dy}{dx} + \frac{P(x)}{y^{n-1}} = Q(x) \label{eq:appendix-ODE-Bernoulli-trans}
\end{gather}
と変形できる。

ここに\dm{u=\frac{1}{y^{n-1}} = y^{1-n}}とおくことで\dm{u'=(1-n)\frac{1}{y^n}\frac{dy}{dx}}
であり、Eq.~\eqref{eq:appendix-ODE-Bernoulli-trans}の両辺を$1-n$倍してから$u$についての微分方程式に書き換えて
\begin{gather}
  u' + (1-n)P(x)u = (1-n)Q(x)
\end{gather}
として一階線形常微分方程式を得る。
